% === Chapter 0: Introduction ===
\section*{Introduction}
\label{sec:t0}
\addcontentsline{toc}{section}{Introduction}

大语言模型（Large Language Model, LLM）和视觉语言模型（Vision-Language Model, VLM）的能力
并非在 pre-training 阶段就已经完全释放。
Pre-training 赋予模型对语言的统计理解和广泛的世界知识，
但一个仅经过 pre-training 的模型本质上只是一个高质量的 next-token predictor ---
它不知道如何遵循人类指令，不能可靠地区分有害输出与安全回答，
也无法在复杂推理任务中展现出结构化的思维过程。
\textbf{Post-training} 正是弥合这一鸿沟的关键阶段：
它将一个 raw language model 转化为一个真正有用的 AI assistant。

\subsection{Post-Training 的定义与重要性}
\label{sec:t0-definition}

本文将 \textbf{post-training} 定义为 pre-training 完成之后、模型部署或 inference 之前，
为提升模型特定能力而进行的所有训练阶段的总称。
这一定义涵盖了以下核心技术范畴：

\begin{itemize}[leftmargin=2em]
  \item \textbf{Supervised Fine-Tuning (SFT)}：
    使用人工标注的 (instruction, response) 数据对模型进行有监督微调，
    使其学会遵循指令并产生格式规范的输出
    \cite{wei2022flan,ouyang2022instructgpt,sanh2022multitask}。

  \item \textbf{Reinforcement Learning from Human Feedback (RLHF)}：
    训练 reward model 捕捉人类偏好，再通过 PPO 等强化学习算法优化模型策略，
    使输出更加符合人类预期
    \cite{christiano2017rlhf,ouyang2022instructgpt,bai2022constitutional}。

  \item \textbf{Direct Preference Optimization (DPO) 及其变体}：
    绕过显式 reward model，直接从偏好数据对中优化策略，
    大幅简化了 alignment 的训练流程
    \cite{rafailov2023dpo,azar2023ipo,ethayarajh2024kto}。

  \item \textbf{其他 post-training 技术}：
    包括 Constitutional AI \cite{bai2022constitutional}、
    process reward model \cite{lightman2023prm}、
    Group Relative Policy Optimization (GRPO) \cite{shao2024grpo}，
    以及面向特定能力（reasoning, safety, tool use, multimodal）的专项训练方法。
\end{itemize}

\noindent
Post-training 与其上下游阶段的边界需要明确区分：

\begin{itemize}[leftmargin=2em]
  \item \textbf{与 pre-training 的区别}：
    Pre-training 在大规模无标注语料上进行，目标是通用语言建模（通常是 next-token prediction 或 masked language modeling）。
    Post-training 则使用相对小规模但高质量的标注数据或偏好信号，目标是赋予模型特定的行为模式。
    两者在数据规模、训练目标和计算开销上存在数量级差异。

  \item \textbf{与 inference-time 技术的区别}：
    Chain-of-thought prompting \cite{wei2022cot}、
    in-context learning \cite{brown2020gpt3}、
    retrieval-augmented generation (RAG) \cite{lewis2020rag}
    等技术在 inference 时通过 prompt 工程或外部知识检索提升模型表现，
    但不修改模型参数。Post-training 则通过梯度更新永久性地改变模型权重。
    当然，两者并非完全正交 ---
    例如 OpenAI o1 \cite{openai2024o1} 将 post-training 阶段习得的推理能力
    与 inference-time 的扩展 compute 策略深度融合。
\end{itemize}

\noindent
Post-training 之所以重要，可以从三个维度理解。
\textbf{第一}，它是模型从"能力"到"可用性"的桥梁 ---
InstructGPT \cite{ouyang2022instructgpt} 以仅 1.3B 参数的 SFT+RLHF 模型在人类评估中胜过
175B 的原始 GPT-3，直接证明了 post-training 可以在不增加模型规模的前提下实现质的飞跃。
\textbf{第二}，它是 alignment 的核心实现手段 ---
从安全对齐到价值观校准，从减少幻觉到提升推理可靠性，
post-training 是将抽象的 alignment 目标转化为可度量训练信号的关键环节。
\textbf{第三}，它是模型差异化的主战场 ---
当 pre-training 架构和数据趋于同质化（大多为 decoder-only Transformer + web-scale corpus），
post-training 策略的差异日益成为决定最终产品质量的关键因素。

\paragraph{本文的目标读者。}
本文面向三类读者：
（1）\textbf{ML 研究者}，希望系统了解 post-training 各子领域的最新进展及其交叉联系；
（2）\textbf{研究生}，需要一份从基础概念到前沿方法的结构化学习路径；
（3）\textbf{工业从业者}，关注 post-training 技术在模型开发中的实际选型与部署策略。
本文假设读者具备 Transformer 架构和 language modeling 的基础知识，
但不要求对每个子领域（如 RLHF、multimodal training、agent tuning）有先验了解。

\paragraph{与既有综述的区别。}
Post-training 相关的综述已有若干代表性工作。
\citet{casper2023rlhfproblems} 系统梳理了 RLHF 的开放问题与根本局限，
但范围限于 reward modeling 和 preference learning，
未涉及 multimodal、agentic 等后续分支。
\citet{wang2023alignsurvey} 以数据收集、训练方法、评估三个维度组织 LLM alignment 技术，
提供了清晰的分类体系，但采用静态分类而非演化视角。
\citet{shen2023alignmentsurvey} 从 inner/outer alignment 的理论框架出发，
覆盖了可解释性和对抗攻击，但未深入各子领域的技术细节。
本文与上述工作在三个方面互补：
（1）\textbf{覆盖范围}更广，将 SFT、preference optimization、safety、reasoning、multimodal、agentic、efficiency 七条线程纳入统一框架；
（2）采用\textbf{演化结构}而非静态分类，追踪每条线程从问题起源到前沿的时间演化；
（3）通过\textbf{跨线程交叉引用}显式记录思想在子领域之间的流动，揭示传统分类方式所忽略的技术迁移。

\subsection{本文的组织方式}
\label{sec:t0-organization}

传统 survey 通常按方法分类（如 SFT 一章、RLHF 一章、DPO 一章），
这种组织方式清晰但静态，无法反映 post-training 领域最核心的特征：
\textbf{不同研究方向之间存在密集的思想流动和技术迁移}。

例如，reward modeling 的思想从 preference optimization（Thread~2）
流向 reasoning verification（Thread~4）中的 process reward model，
再通过 GRPO 回流到 preference optimization；
SFT 的 instruction tuning pipeline（Thread~1）
被 multimodal post-training（Thread~5）和 tool use（Thread~6）直接借鉴和扩展；
efficiency 技术如 LoRA（Thread~7）使得前述所有线程的方法能够以更低成本落地。

为了捕捉这种动态结构，本文采用 \textbf{Hybrid Problem-Evolution Graph} 作为组织框架。
其核心思想如下：

\begin{enumerate}[leftmargin=2em]
  \item \textbf{七条问题驱动的演化线程}：
    本文将 post-training 领域划分为七条相对独立的问题线程（thread），
    每条线程围绕一个核心问题展开。
    线程内部按时间顺序追踪
    ``Problem Origin $\rightarrow$ Foundation $\rightarrow$ Limitations \& Branching $\rightarrow$ Convergence $\rightarrow$ Open Problems''
    的演化链条。

  \item \textbf{跨线程交叉引用}：
    当某条线程中的技术或思想被另一条线程借鉴时，
    本文使用 \crossref{1}{sec:t1} 样式的交叉引用进行标注，
    使读者可以追踪思想的来源和去向。

  \item \textbf{全景视图}：
    \cref{fig:overview} 以 TikZ 图的形式呈现了七条线程的全景视图，
    其中水平泳道对应线程，x 轴对应时间线，
    节点表示 landmark 工作，虚线箭头表示跨线程的思想流动。
\end{enumerate}

\noindent
这种组织方式的优势在于：
（1）在线程内部保持了清晰的时间线和因果关系；
（2）通过跨线程引用揭示了被传统分类方式割裂的思想联系；
（3）全景图为读者提供了在任意位置开始阅读的入口。

\subsection{七条演化线程概述}
\label{sec:t0-threads}

\paragraph{Thread 1: Instruction Following \& SFT.}
核心问题：\textit{如何让 pre-trained language model 学会理解并遵循人类指令？}
FLAN \cite{wei2022flan} 和 T0 \cite{sanh2022multitask} 首先证明了
multi-task instruction tuning 可以显著提升模型的 zero-shot 泛化能力。
InstructGPT \cite{ouyang2022instructgpt} 将 SFT 作为 RLHF pipeline 的第一阶段，
确立了"SFT $\rightarrow$ reward model $\rightarrow$ PPO"的经典三阶段训练范式。
Self-Instruct \cite{wang2023selfinstruct} 开创了利用模型自身生成 instruction 数据的方法，
极大降低了标注成本并催生了 Alpaca \cite{taori2023alpaca} 等开源 instruction-tuned 模型。
LIMA \cite{zhou2023lima} 则以仅 1,000 条精心筛选的样本达到了极具竞争力的性能，
提出了"Superficial Alignment Hypothesis" --- 即 alignment 所需的知识已在 pre-training 中习得，
SFT 的作用仅在于激活正确的输出格式和风格。
\crossref{1}{sec:t1}

\paragraph{Thread 2: Reward Modeling \& Preference Optimization.}
核心问题：\textit{如何将模糊的人类偏好转化为可优化的训练信号？}
RLHF 的基础框架可追溯至
Christiano et al.\ \cite{christiano2017rlhf} 在 Atari 和 MuJoCo 任务上的工作，
随后被 Stiennon et al.\ \cite{stiennon2020summarize} 应用于文本摘要，
并在 InstructGPT \cite{ouyang2022instructgpt} 中全面应用于对话模型的训练。
然而 RLHF 的训练流程复杂且不稳定，涉及四个独立模型的协同训练。
DPO \cite{rafailov2023dpo} 通过数学推导证明 reward model 可以被隐式参数化，
从而将 RLHF 转化为一个简单的分类损失，大幅降低了实现门槛。
此后 IPO \cite{azar2023ipo}、KTO \cite{ethayarajh2024kto}、
ORPO \cite{hong2024orpo} 等变体从不同角度改进了 DPO 的局限性。
GRPO \cite{shao2024grpo} 则将 group-relative reward 引入 policy optimization，
在 reasoning 任务上取得了突破性进展。
2025 年，RLVR（Reinforcement Learning with Verifiable Rewards）范式的兴起
标志着新的分水岭：DAPO \cite{yu2025dapo} 等方法完全抛弃 reward model，
直接使用规则验证器作为 reward 信号，将 GRPO 在 AIME 上的表现从 30\% 提升至 50\%。
\crossref{2}{sec:t2}

\paragraph{Thread 3: Safety \& Alignment.}
核心问题：\textit{如何确保模型在遵循指令的同时拒绝有害请求，且不过度保守？}
Constitutional AI (CAI) \cite{bai2022constitutional} 提出让模型根据一组预定义原则（constitution）
进行自我批评和修改，减少了 safety alignment 对人工标注有害数据的依赖。
Safe RLHF \cite{dai2024saferlhf} 将 helpfulness 和 harmlessness 解耦为两个独立目标，
通过带约束的优化框架在 Pareto 前沿上寻求最优平衡。
与此同时，jailbreaking 攻击的发展 ---
从简单的 prompt injection 到 GCG \cite{zou2023gcg} 等基于梯度的自动攻击 ---
形成了一场持续的攻防军备竞赛，推动 safety 研究不断迭代。
2025--2026 年，推理模型的出现带来了新的安全维度：
研究表明 reasoning model 可被用作自主 jailbreak agent，
在零人工干预下达到 97\% 的攻击成功率 \cite{hagendorff2026jailbreakagents}，
推理链（chain-of-thought）的安全性成为新的核心挑战。
\crossref{3}{sec:t3}

\paragraph{Thread 4: Reasoning \& Verification.}
核心问题：\textit{如何让模型进行可靠的多步推理，并对中间步骤进行验证？}
Chain-of-Thought (CoT) prompting \cite{wei2022cot} 证明了
在 prompt 中加入推理过程可以显著提升模型在数学和逻辑任务上的表现，
但这种能力的上限受限于 inference-time 的 prompt 设计。
Process Reward Model (PRM) \cite{lightman2023prm} 将验证粒度从最终答案细化到每一步推理，
为 reward-guided search 提供了更精细的信号。
OpenAI o1 \cite{openai2024o1} 将 reasoning 能力深度融入 post-training，
通过大规模 reinforcement learning 训练模型在 inference 时自主展开长链推理，
在数学竞赛和编程任务上达到了专家级水平。
DeepSeek-R1 \cite{deepseek2025r1} 进一步探索了
纯 RL（无 SFT cold start）训练推理能力的可能性，
揭示了 reasoning 能力可以从 RL 中自发涌现的惊人现象。
2025--2026 年，thinking models 成为默认范式：
OpenAI o3 \cite{openai2025o3} 在 AIME 2025 上达到 88.9\%，
DeepSeek-V3.2 \cite{deepseek2025v32} 在 IMO 级别数学竞赛中达到金牌水平，
Qwen3 \cite{yang2025qwen3} 实现了 thinking/non-thinking 模式的无缝融合。
\crossref{4}{sec:t4}

\paragraph{Thread 5: Multimodal Post-training.}
核心问题：\textit{如何将 text-only 的 post-training 技术扩展到视觉-语言联合建模？}
LLaVA \cite{liu2023llava} 将 visual instruction tuning 引入 VLM 训练，
通过两阶段训练（visual-text alignment + instruction tuning）
构建了高效的多模态 instruction-following 模型。
随后，RLHF 技术被引入 VLM 以减少 multimodal hallucination ---
即模型生成与图像内容不一致的描述
\cite{sun2023rlhfv,yu2024rlhfv}。
Hallucination reduction 成为 multimodal post-training 的核心挑战之一，
催生了 POPE \cite{li2023pope}、
CHAIR 等专项评估基准和针对性训练方法。
2025 年，RL for VLM reasoning 成为新的前沿方向：
R1-VL \cite{zhang2025r1vl}、MM-Eureka \cite{meng2025mmeureka} 等工作
将 RLVR 范式从 text-only 推理扩展到多模态推理，
标志着 VLM post-training 从 SFT+DPO 向 SFT+RL 的范式转型。
\crossref{5}{sec:t5}

\paragraph{Thread 6: Tool Use \& Agentic Capabilities.}
核心问题：\textit{如何通过 post-training 让模型学会调用外部工具并完成复杂任务？}
Toolformer \cite{schick2023toolformer} 证明了模型可以通过 self-supervised learning
学会在生成过程中插入 API 调用。
OpenAI 在 GPT-3.5/4 中引入的 function calling 机制
将 tool use 标准化为结构化的 JSON schema 接口，
使其成为生产级应用的基础设施。
Agent tuning \cite{zeng2023agenttuning} 进一步将训练目标从单次工具调用
扩展到多步骤任务规划和执行，为构建 autonomous agent 奠定了基础。
这条线程与 Thread~1（instruction following 提供了基础能力）
和 Thread~4（reasoning 提供了规划能力）存在深度交叉。
2024--2025 年，agent 能力从 API 调用进化为直接操作计算机 GUI：
Claude Computer Use \cite{anthropic2024computeruse}、OpenAI Operator \cite{openai2025operator}
展示了 AI agent 在真实桌面环境中完成复杂任务的能力，
MCP \cite{anthropic2024mcp} 等标准化协议正在将 agent 从实验原型推向工业部署。
\crossref{6}{sec:t6}

\paragraph{Thread 7: Efficiency \& Data Engineering.}
核心问题：\textit{如何以更低的计算成本和更少的数据实现同等或更优的 post-training 效果？}
LoRA \cite{hu2022lora} 通过低秩分解将 fine-tuning 的可训练参数量降低至原始模型的 0.1\%--1\%，
使得在消费级 GPU 上进行 post-training 成为可能。
QLoRA \cite{dettmers2023qlora} 进一步结合量化技术，
在 48GB 显存内实现了 65B 模型的 fine-tuning。
在数据工程维度，
data selection 方法（如 DEITA \cite{liu2024deita}、
Instruction Mining \cite{cao2024instructionmining}）致力于从大规模 instruction 数据中筛选最有价值的子集；
synthetic data generation 则利用强模型（如 GPT-4）生成训练数据，
催生了 Alpaca \cite{taori2023alpaca}、
WizardLM \cite{xu2024wizardlm}、
Evol-Instruct 等方法论。
\crossref{7}{sec:t7}

\subsection{跨线程交叉引用}
\label{sec:t0-crossref}

七条演化线程并非独立发展，它们之间存在密集的思想流动和技术迁移。
\cref{fig:overview} 以全景视图的方式展示了这些跨线程联系。
这里我们概述最重要的几条 idea flow：

\begin{itemize}[leftmargin=2em]
  \item \textbf{T1 $\rightarrow$ T2}（SFT 为 RLHF 提供 warm start）：
    InstructGPT 确立的 ``SFT $\rightarrow$ RM $\rightarrow$ PPO'' pipeline
    将 instruction tuning 作为 preference optimization 的前置阶段，
    这一范式被几乎所有后续工作继承。

  \item \textbf{T1 $\rightarrow$ T3}（Instruction following 暴露 safety 风险）：
    T1 中 SFT 赋予模型的指令遵循能力同时打开了安全漏洞——
    一个忠实遵循指令的模型也会忠实执行恶意指令。
    这一内在张力正是 T3 的问题起源，
    驱动了 Constitutional AI 等安全对齐方法的发展。

  \item \textbf{T2 $\rightarrow$ T3}（Preference optimization 塑造 safety）：
    RLHF 和 DPO 提供了将人类安全偏好转化为训练信号的通用框架，
    Constitutional AI 和 Safe RLHF 在此基础上发展出面向 safety 的专门化方法。

  \item \textbf{T2 $\leftrightarrow$ T4}（Reward modeling 与 reasoning 的双向流动）：
    Reward model 的思想从 T2 流入 T4，催生了 process reward model；
    反过来，T4 中发展的 verifiable reward（基于 ground truth 答案的自动验证）
    启发了 T2 中 GRPO 等方法，实现了无需人工标注的 preference optimization。

  \item \textbf{T1 $\rightarrow$ T5, T6}（Instruction tuning 的跨模态迁移）：
    T1 中发展成熟的 instruction tuning pipeline 被 T5（LLaVA 的 visual instruction tuning）
    和 T6（Toolformer 的 tool use instruction tuning）直接借鉴，
    证明了这一范式的高度可迁移性。

  \item \textbf{T7 $\rightarrow$ T1--T6}（Efficiency 技术的普惠效应）：
    LoRA、QLoRA 等参数高效训练方法使得前述六条线程中的 post-training 技术
    能够在学术界和开源社区广泛复现和迭代，显著加速了整个领域的发展。

  \item \textbf{T4 $\rightarrow$ T6}（Reasoning 赋能 agent planning）：
    T4 中发展的长链推理能力为 T6 中的 multi-step agent planning 提供了核心支撑，
    o1 和 DeepSeek-R1 的 reasoning 能力直接提升了 agent 在复杂任务中的表现。

  \item \textbf{T2 $\rightarrow$ T4 $\rightarrow$ T5}（RLVR 的跨线程迁移）：
    2025 年，RLVR 范式在 T2 中被系统化（DAPO \cite{yu2025dapo}），
    在 T4 中被大规模应用于 reasoning 训练，
    随后迅速迁移到 T5（R1-VL \cite{zhang2025r1vl}、MM-Eureka \cite{meng2025mmeureka}），
    形成了清晰的 T2 $\rightarrow$ T4 $\rightarrow$ T5 技术迁移路径。
\end{itemize}

% === Overview: Hybrid Problem-Evolution Graph ===
% 7 horizontal lanes (one per thread), shared x-axis 2017--2026,
% landmark papers as nodes, dashed cross-thread arrows.

\begin{figure}[htbp]
\centering
\resizebox{\textwidth}{!}{%
\begin{tikzpicture}[
    >=Stealth,
    node distance=0.6cm,
    paper/.style={
        rounded corners=2pt,
        font=\footnotesize\bfseries,
        inner sep=4pt,
        minimum height=0.65cm,
        text=white,
        drop shadow={shadow xshift=0.5pt, shadow yshift=-0.5pt, opacity=0.3},
    },
    crossarrow/.style={
        ->,
        dashed,
        very thick,
        shorten >=2pt,
        shorten <=2pt,
    },
    arrowlabel/.style={
        font=\tiny\itshape,
        fill=white,
        fill opacity=0.85,
        text opacity=1,
        inner sep=1pt,
        rounded corners=1pt,
    },
    lanelabel/.style={
        font=\small\bfseries,
        anchor=east,
        text width=4.2cm,
        align=right,
    },
]

% --- Geometry ---
\def\laneheight{1.7}   % vertical spacing between lanes
\def\xmin{0}
\def\xmax{20}          % total width
% Year mapping: 2017 = 0, 2026 = 20  =>  1 year = 2 units
\newcommand{\yearx}[1]{\fpeval{(#1-2017)*2}}

% --- Year axis ---
\foreach \y in {2017,2018,2019,2020,2021,2022,2023,2024,2025,2026} {
    \draw[gray!40] (\yearx{\y}, 0.5*\laneheight) -- (\yearx{\y}, -6.5*\laneheight);
    \node[font=\footnotesize, gray!80, anchor=south] at (\yearx{\y}, 0.7*\laneheight) {\y};
}

% --- Lane backgrounds ---
\foreach \i/\col in {0/thread1,1/thread2,2/thread3,3/thread4,4/thread5,5/thread6,6/thread7} {
    \fill[\col, opacity=0.06]
        (\xmin, -\i*\laneheight+0.45*\laneheight)
        rectangle
        (\xmax, -\i*\laneheight-0.45*\laneheight);
}

% =====================================================================
% Thread 1: Instruction Following & SFT (lane y=0)
% =====================================================================
\def\laneA{0}
\node[lanelabel, text=thread1] at (-0.3, \laneA*\laneheight)
    {T1: Instruction\\Following \& SFT};

\node[paper, fill=thread1] (flan)
    at (\yearx{2021}+0.5, \laneA*\laneheight) {FLAN};
\node[paper, fill=thread1] (instructgpt-sft)
    at (\yearx{2022}+0.8, \laneA*\laneheight) {InstructGPT};
\node[paper, fill=thread1] (selfinstruct)
    at (\yearx{2023}+1.2, \laneA*\laneheight) {Self-Inst.};
\node[paper, fill=thread1] (lima)
    at (\yearx{2024}+0.5, \laneA*\laneheight) {LIMA};

% =====================================================================
% Thread 2: Reward Modeling & Preference Optimization (lane y=-1)
% =====================================================================
\def\laneB{-1}
\node[lanelabel, text=thread2] at (-0.3, \laneB*\laneheight)
    {T2: Reward Modeling \&\\Preference Opt.};

\node[paper, fill=thread2] (rlhf-og)
    at (\yearx{2017}+1.0, \laneB*\laneheight) {RLHF};
\node[paper, fill=thread2] (instructgpt-rlhf)
    at (\yearx{2022}+0.8, \laneB*\laneheight) {InstructGPT};
\node[paper, fill=thread2] (dpo)
    at (\yearx{2023}+0.6, \laneB*\laneheight) {DPO};
\node[paper, fill=thread2] (grpo)
    at (\yearx{2024}+1.0, \laneB*\laneheight) {GRPO};
\node[paper, fill=thread2] (dapo)
    at (\yearx{2025}+1.0, \laneB*\laneheight) {DAPO/RLVR};

% =====================================================================
% Thread 3: Safety & Alignment (lane y=-2)
% =====================================================================
\def\laneC{-2}
\node[lanelabel, text=thread3] at (-0.3, \laneC*\laneheight)
    {T3: Safety \&\\Alignment};

\node[paper, fill=thread3] (cai)
    at (\yearx{2022}+0.8, \laneC*\laneheight) {Const.\ AI};
\node[paper, fill=thread3] (saferlhf)
    at (\yearx{2023}+1.2, \laneC*\laneheight) {Safe RLHF};
\node[paper, fill=thread3] (jailbreak)
    at (\yearx{2024}+1.2, \laneC*\laneheight) {Jailbreak};

% =====================================================================
% Thread 4: Reasoning & Verification (lane y=-3)
% =====================================================================
\def\laneD{-3}
\node[lanelabel, text=thread4] at (-0.3, \laneD*\laneheight)
    {T4: Reasoning \&\\Verification};

\node[paper, fill=thread4] (cot)
    at (\yearx{2022}+0.3, \laneD*\laneheight) {CoT};
\node[paper, fill=thread4] (prm)
    at (\yearx{2023}+0.5, \laneD*\laneheight) {PRM};
\node[paper, fill=thread4] (oone)
    at (\yearx{2024}+0.0, \laneD*\laneheight) {o1};
\node[paper, fill=thread4] (deepseekr1)
    at (\yearx{2025}-0.3, \laneD*\laneheight) {R1};
\node[paper, fill=thread4] (othree)
    at (\yearx{2025}+0.8, \laneD*\laneheight) {o3};
\node[paper, fill=thread4] (v32)
    at (\yearx{2026}+0.0, \laneD*\laneheight) {V3.2};

% =====================================================================
% Thread 5: Multimodal Post-training (lane y=-4)
% =====================================================================
\def\laneE{-4}
\node[lanelabel, text=thread5] at (-0.3, \laneE*\laneheight)
    {T5: Multimodal\\Post-training};

\node[paper, fill=thread5] (llava)
    at (\yearx{2023}-0.3, \laneE*\laneheight) {LLaVA};
\node[paper, fill=thread5] (vlmrlhf)
    at (\yearx{2024}+0.0, \laneE*\laneheight) {VLM-RLHF};
\node[paper, fill=thread5] (halluc)
    at (\yearx{2025}+0.3, \laneE*\laneheight) {Anti-Halluc.};

% =====================================================================
% Thread 6: Tool Use & Agentic (lane y=-5)
% =====================================================================
\def\laneF{-5}
\node[lanelabel, text=thread6] at (-0.3, \laneF*\laneheight)
    {T6: Tool Use \&\\Agentic};

\node[paper, fill=thread6] (toolformer)
    at (\yearx{2022}+1.5, \laneF*\laneheight) {Toolformer};
\node[paper, fill=thread6] (funccall)
    at (\yearx{2023}+1.5, \laneF*\laneheight) {Func.\ Call};
\node[paper, fill=thread6] (agenttune)
    at (\yearx{2024}+1.5, \laneF*\laneheight) {AgentTune};
\node[paper, fill=thread6] (mcp)
    at (\yearx{2025}+1.3, \laneF*\laneheight) {MCP/Agents};

% =====================================================================
% Thread 7: Efficiency & Data Engineering (lane y=-6)
% =====================================================================
\def\laneG{-6}
\node[lanelabel, text=thread7] at (-0.3, \laneG*\laneheight)
    {T7: Efficiency \&\\Data Engineering};

\node[paper, fill=thread7] (lora)
    at (\yearx{2021}+1.0, \laneG*\laneheight) {LoRA};
\node[paper, fill=thread7] (datasel)
    at (\yearx{2023}+0.8, \laneG*\laneheight) {Data Selection};
\node[paper, fill=thread7] (synthdata)
    at (\yearx{2024}+0.5, \laneG*\laneheight) {Synthetic Data};

% =====================================================================
% Cross-thread arrows (dashed, idea flow)
% =====================================================================

% T1 -> T2: InstructGPT-SFT feeds into InstructGPT-RLHF
\draw[crossarrow, thread1] (instructgpt-sft.south) -- (instructgpt-rlhf.north);

% T2 -> T3: RLHF ideas feed into Constitutional AI
\draw[crossarrow, thread2] (instructgpt-rlhf.south) -- (cai.north)
    node[arrowlabel, midway, right=1pt] {reward};

% T2 -> T4: Reward modeling feeds reasoning verification (PRM)
\draw[crossarrow, thread2] (dpo.south)
    to[out=-90, in=90] (prm.north);

% T4 -> T2: Reasoning verification inspires GRPO
\draw[crossarrow, thread4] (prm.east)
    to[out=0, in=180] (grpo.south west)
    node[arrowlabel, midway, above=1pt] {verification};

% T1 -> T5: SFT pipeline adopted by LLaVA
\draw[crossarrow, thread1] (selfinstruct.south)
    to[out=-90, in=90] (llava.north)
    node[arrowlabel, near start, right=1pt] {SFT pipeline};

% T2 -> T5: RLHF adopted for VLM alignment
\draw[crossarrow, thread2] (dpo.south)
    to[out=-100, in=90] (vlmrlhf.north)
    node[arrowlabel, near end, left=1pt] {preference};

% T1 -> T6: Instruction tuning enables tool use
\draw[crossarrow, thread1] (selfinstruct.south)
    to[out=-90, in=90] (toolformer.north);

% T7 -> T1: LoRA enables efficient SFT
\draw[crossarrow, thread7] (lora.north)
    to[out=90, in=-90] (flan.south)
    node[arrowlabel, near start, right=1pt] {LoRA};

% T7 -> T2: Efficient training enables broader RLHF
\draw[crossarrow, thread7] (datasel.north)
    to[out=90, in=-90] (grpo.south);

% T3 -> T4: Safety constraints shape reasoning alignment
\draw[crossarrow, thread3] (saferlhf.south)
    to[out=-90, in=90] (oone.north)
    node[arrowlabel, midway, right=1pt] {safety};

% T4 -> T6: Reasoning enables agent planning
\draw[crossarrow, thread4] (oone.south)
    to[out=-90, in=90] (agenttune.north)
    node[arrowlabel, midway, right=1pt] {reasoning};

% T2 -> T4: RLVR (DAPO) feeds reasoning training
\draw[crossarrow, thread2] (dapo.south)
    to[out=-90, in=90] (othree.north)
    node[arrowlabel, near start, right=1pt] {RLVR};

% T4 -> T5: Reasoning RL migrates to multimodal (R1-VL, MM-Eureka)
\draw[crossarrow, thread4] (deepseekr1.south)
    to[out=-90, in=90] (halluc.north east);

% T4 -> T6: Thinking models empower agent planning
\draw[crossarrow, thread4] (othree.south)
    to[out=-90, in=90] (mcp.north);

\end{tikzpicture}%
}% end resizebox
\caption{Post-training 领域的 Hybrid Problem-Evolution Graph 全景视图。
七条水平泳道分别对应七条演化线程，
x 轴为时间线（2017--2026），
节点为各线程内的 landmark 工作，
虚线箭头表示跨线程的思想流动与技术迁移。}
\label{fig:overview}
\end{figure}


\noindent
在接下来的章节中，我们将沿着每条线程深入分析 landmark 论文的技术细节，
同时通过跨线程引用保持全局视角。读者可以按照线程顺序（T1--T7）线性阅读，
也可以从任意感兴趣的线程入手，通过交叉引用跳转到相关线程。
